\section{Electromagnetically Induced Transparency Memory}
\label{sec:eit_memory}

EIT storage is enabled by a $\Lambda$-type atomic system in which two long-lived ground states, denoted by $\ket{g}$ and $\ket{s}$, are optically coupled to a common excited state $\ket{e}$. The weak signal field couples the transition $\ket{g}\leftrightarrow\ket{e}$, while a stronger classical control field couples $\ket{s}\leftrightarrow\ket{e}$. Under suitable resonance conditions, the control field modifies the optical response of the ensemble so that resonant absorption is suppressed and the optical excitation can be mapped coherently onto a ground-state spin coherence~\cite{Apgteitiav_FinkelsteinEtAl_2023,Psliwavueit_DeRoseEtAl_2023a,}.

In the absence of the control field, the signal field addresses an ordinary resonant optical transition. Atoms initially prepared in $\ket{g}$ may absorb signal photons and be excited to $\ket{e}$. Since the excited state is short-lived, this excitation decays by spontaneous emission and removes the photon in a random optical mode. An optically dense ensemble is therefore opaque to a resonant signal field unless the optical response is modified by the control field~\cite{Apgteitiav_FinkelsteinEtAl_2023,Psliwavueit_DeRoseEtAl_2023a}.

When the signal and control fields satisfy two-photon resonance, the atom can be prepared in a coherent superposition of the two ground states with no excited-state component. In the standard $\Lambda$-system description, the dark state is written as~\cite{Apgteitiav_FinkelsteinEtAl_2023,Psliwavueit_DeRoseEtAl_2023a}
\begin{equation}
    \ket{D}
    =
    \frac{
    \Omega_{\mathrm{c}}\ket{g}
    -
    \Omega_{\mathrm{sig}}\ket{s}
    }{
    \sqrt{|\Omega_{\mathrm{c}}|^2+|\Omega_{\mathrm{sig}}|^2}
    },
    \label{eq:eit_dark_state}
\end{equation}
where $\Omega_{\mathrm{sig}}$ and $\Omega_{\mathrm{c}}$ are the Rabi frequencies of the weak signal field on $\ket{g}\leftrightarrow\ket{e}$  and the control field on $\ket{s}\leftrightarrow\ket{e}$ respectively. Absorption is suppressed because the two optical excitation pathways interfere destructively at $\ket{e}$. The ensemble can therefore remain transparent to a resonant signal field even though the atoms occupy a ground state that would absorb in the absence of the control field~\cite{Apgteitiav_FinkelsteinEtAl_2023,Psliwavueit_DeRoseEtAl_2023a}.

The destructive interference also reshapes the optical susceptibility. In the weak-signal limit, the signal field probes the linear response of the atomic medium, while the control field sets the width and strength of the EIT resonance. A commonly used expression for the signal susceptibility is~\cite{Apgteitiav_FinkelsteinEtAl_2023,}
\begin{equation}
    \chi_{\mathrm{sig}}(\Delta,\delta)
    =
        i\frac{\mathcal{N}|\mu_{ge}|^2}{\hbar\epsilon_0}
    \frac{\Gamma_{gs}}
    {\Gamma_{gs}\Gamma_{ge}+|\Omega_{\mathrm{c}}|^2},
    \label{eq:eit_susceptibility}
\end{equation}

where $\mathcal{N}$ is the atomic number density, $\mu_{ge}$ is the electric-dipole matrix element of the signal transition, and $\epsilon_0$ is the vacuum permittivity. The optical detuning, defined as $\Delta=\omega_{\mathrm{sig}}-\omega_{eg}$, where $\omega_{\mathrm{sig}}$ is the signal-field angular frequency and $\omega_{eg}$ is the resonance frequency of the $\ket{g}\leftrightarrow\ket{e}$ transition. The two-photon detuning $\delta=(\omega_{\mathrm{sig}}-\omega_{\mathrm{c}})-\omega_{sg}$, where $\omega_{\mathrm{c}}$ is the control-field angular frequency and $\omega_{sg}$ is the ground-state splitting between $\ket{g}$ and $\ket{s}$. The complex response factors $\Gamma_{ge}=\gamma_{ge}-i\Delta$ and $\Gamma_{gs}=\gamma_{gs}-i\delta$ combine coherent detuning with irreversible loss of coherence. Here, $\gamma_{ge}$ is the decay rate of the optical coherence associated with $\ket{g}\leftrightarrow\ket{e}$, while $\gamma_{gs}$ is the decoherence rate of the stored ground-state coherence between $\ket{g}$ and $\ket{s}$.

At exact two-photon resonance, \(\delta=0\). If the ground-state coherence is long lived, so that \(\gamma_{gs}\) is negligible, then \(\Gamma_{gs}\rightarrow0\). In Eq.~\eqref{eq:eit_susceptibility}, the numerator therefore tends to zero, and the signal susceptibility vanishes at the centre of the EIT resonance. Physically, the medium no longer absorbs the resonant signal field because the excitation pathway through \(\ket{e}\) is cancelled by destructive interference. Away from the exact resonance, however, the cancellation is incomplete. The imaginary part of \(\chi_{\mathrm{sig}}\), which describes absorption, remains strongly suppressed inside the transparency window, while the real part changes rapidly with detuning. This steep variation of the real part is the origin of the large group index and the slow propagation of the signal pulse~\cite{Apgteitiav_FinkelsteinEtAl_2023,}.

This steep dispersion inside the EIT window reduces the group velocity of the signal pulse. For a dispersive medium, the scalar wavenumber of the signal field is~\cite{Apgteitiav_FinkelsteinEtAl_2023,}
\begin{equation}
    k(\omega)
    =
    \frac{n(\omega)\omega}{c},
    \label{eq:eit_dispersion_relation}
\end{equation}
where $\omega$ is the optical angular frequency, $n(\omega)$ is the refractive index, and $c$ is the speed of light in vacuum. The group velocity is then~\cite{Apgteitiav_FinkelsteinEtAl_2023,Psliwavueit_DeRoseEtAl_2023a}
\begin{equation}
    v_{\mathrm{g}}
    =
    \left(\frac{\partial k}{\partial\omega}\right)^{-1}
    =
    \frac{c}
    {n+\omega\,\frac{\partial n}{\partial\omega}},
    \label{eq:eit_group_velocity}
\end{equation}
where $v_{\mathrm{g}}$ is the velocity of the pulse envelope. Inside the EIT window, absorption is small and $\partial n/\partial\omega$ can be large, so the denominator in Eq.~\eqref{eq:eit_group_velocity} may greatly exceed unity. The pulse then propagates through the ensemble with a group velocity far below $c$~\cite{Apgteitiav_FinkelsteinEtAl_2023,Psliwavueit_DeRoseEtAl_2023a,}.

The reduced group velocity compresses the signal pulse inside the ensemble, which allows the optical mode to be spatially contained before the control field is extinguished. The front of the pulse is slowed after entering the EIT medium, while the remaining part of the pulse continues to enter the cell. For sufficient optical depth and an appropriate control-field strength, the full pulse can be placed inside the ensemble. Only the optical excitation inside the medium can be converted into atomic coherence during storage~\cite{QmfpDp_FleischhauerEtAl_2002,Psliwavueit_DeRoseEtAl_2023a,SoLiAV_PhillipsEtAl_2001,}.

In the dark-state polariton treatment of Fleischhauer and Lukin, storage is described as the adiabatic rotation of a joint light--matter excitation rather than as ordinary resonant absorption~\cite{QmfpDp_FleischhauerEtAl_2002}. The polariton field is written as
\begin{equation}
    \hat{\Psi}(z,t)
    =
    \cos\theta(t)\,\hat{\mathcal{E}}(z,t)
    -
    \sin\theta(t)\,\hat{S}(z,t),
    \label{eq:dark_state_polariton}
\end{equation}

where $\hat{\Psi}(z,t)$ is the dark-state polariton field, $\hat{\mathcal{E}}(z,t)$ is the slowly varying quantum field operator of the signal light, and $\hat{S}(z,t)$ is the collective ground-state spin coherence of the ensemble. The polariton is a coherent superposition of a propagating optical excitation and a stationary atomic coherence. Its photonic weight is $\cos^2\theta(t)$, while its atomic weight is $\sin^2\theta(t)$. The mixing angle $\theta(t)$ is therefore the parameter that determines where the quantum excitation is stored at a given time.

Following the dark-state polariton treatment of Fleischhauer and Lukin~\cite{QmfpDp_FleischhauerEtAl_2002}, the mixing angle is set by the ratio between the collective atom--light coupling and the control-field Rabi frequency,
\begin{equation}
    \tan^2\theta(t)
    =
    \frac{g^2N}{|\Omega_{\mathrm{c}}(t)|^2}.
    \label{eq:eit_mixing_angle}
\end{equation}
Here, $g$ is the single-atom coupling strength to the signal mode, $N$ is the effective number of atoms coupled to that mode, and $g\sqrt{N}$ is the collectively enhanced coupling between the ensemble and the signal field. The control Rabi frequency $\Omega_{\mathrm{c}}(t)$ controls the coupling between the optical polarization and the spin coherence. A large $|\Omega_{\mathrm{c}}|$ makes $g^2N/|\Omega_{\mathrm{c}}|^2$ small, so $\theta\approx0$ and $\hat{\Psi}$ is predominantly optical. As the control field is reduced, $\theta$ increases and the same excitation acquires a larger spin-wave component. In the limit $|\Omega_{\mathrm{c}}|\rightarrow0$, $\theta\rightarrow\pi/2$, the photonic weight $\cos^2\theta$ vanishes, and the excitation is mapped onto the collective ground-state coherence $\hat{S}(z,t)$~\cite{QmfpDp_FleischhauerEtAl_2002,}.

Storage is obtained when the control field is reduced adiabatically while the system remains in the dark state. The optical part of the polariton is then converted into the ground-state coherence between $\ket{g}$ and $\ket{s}$, while population of $\ket{e}$ remains suppressed. The stored excitation can survive longer than an optical polarization because it is carried by long-lived ground states rather than by the short-lived excited state~\cite{QmfpDp_FleischhauerEtAl_2002,SoLiAV_PhillipsEtAl_2001,}.

Retrieval is produced by applying the control field again. The mixing angle is rotated back toward the photonic limit, the stored spin coherence is converted into an optical polarization, and the signal field is regenerated. If decoherence and mode distortion during storage are negligible, the retrieved pulse preserves the quantum state, phase, and temporal mode of the input. In practice, retrieval is limited by the lifetime and spatial structure of the stored coherence, as well as by optical depth, bandwidth, and noise~\cite{QmfpDp_FleischhauerEtAl_2002,Apgteitiav_FinkelsteinEtAl_2023}.

The full time-dependent EIT writing dynamics are not solved explicitly in the retrieval model presented later. EIT is used to justify the initial condition after storage, where the travelling optical excitation has been converted into a collective ground-state coherence. The relevant object is then the spin wave left by the storage process, whose spatial envelope and phase determine the optical field that can be retrieved~\cite{QmfpDp_FleischhauerEtAl_2002,}.


